\chapter{Machine Learning}
\label{ch:ml}

Most of this book computes an answer from a rule. You give Mathilda a polynomial
and it factors it; you give it a differential equation and it solves it. Machine
learning runs the other way. You give it examples of the answer, and it builds the
rule. A \emph{classifier}\index{classifier} learns to name the class of a new point
from points whose classes are known. A \emph{predictor} learns a numerical response
in the same way. A \emph{clustering} finds groups in points that carry no labels at
all, and a \emph{density estimate} learns the distribution the points came from.
This chapter covers all four, with the mathematics of each method and the check
that confirms it at the prompt.

Mathilda's machine-learning builtins live in \texttt{src/ml/}. \B{Classify} offers
five methods: nearest neighbours, naive Bayes, logistic regression, decision trees,
and random forests. \B{Predict} offers least-squares regression and nearest
neighbours. \B{FindClusters} has ten clustering methods, \B{LearnDistribution}
fits Gaussians, mixtures, and categorical tables, and \B{PrincipalComponents},
\B{DimensionReduce}, and \B{DimensionReduction} reduce dimension. The set is
deliberately small next to Mathematica's. Every method is a textbook algorithm
with its decisions stated in the source, and every fitted model is an ordinary
expression that you can take apart with \B{Part}. The chapter uses that openness
throughout: when a method claims a decision boundary at \(4.5\), the transcript
reads the coefficients back and checks it.

The chapter has three parts. The first builds classifiers on one running example
and explains each method. The second covers evaluation, regression, clustering,
mixtures, and principal components. The last section is a set of short featured
examples on real data: Fisher's irises, the 1974 Motor Trend road tests, spam
messages, and handwritten digits.

\section{A first classifier}
\label{sec:ml-first}

The running example is a set of ninety points in the plane, thirty from each of
three overlapping Gaussian clouds. Each point carries a class, \mcode{"A"},
\mcode{"B"}, or \mcode{"C"}. Here are the points coloured by class. The clouds are
centred at \((0,0)\), \((3,3)\), and \((0,4)\), each with standard deviation
\(0.8\), so the second and third clouds overlap noticeably.

\plotcell[0.6\linewidth]{machine-learning/blobs-data}

The helper \mcode{blob} draws its coordinates from \B{RandomVariate}, which uses
the same stream as \B{RandomReal}, so \B{SeedRandom} makes the whole session
reproducible. Every session in this chapter starts from a seed, and running it
twice gives the same transcript.

\subsection{Training data and the fitted model}
\label{ssec:ml-data}

\mcode{Classify} takes its training data as a list of rules
\(\text{features} \to \text{class}\). The features are a list of numbers, or a
single number when there is only one feature. \B{Thread} turns a list of points
and a class into such rules:

\begin{codepairs}
\pair{machine-learning/first/1}{A seeded generator of Gaussian clouds.}
\pair{machine-learning/first/2}{Three clouds of thirty points.}
\pair{machine-learning/first/3}{Each point becomes a rule \(\text{point} \to
\text{class}\).}
\pair{machine-learning/first/4}{Ninety training examples.}
\pair{machine-learning/first/5}{The first two examples.}
\pair{machine-learning/first/6}{\mcode{Classify} returns a \B{ClassifierFunction}.
It prints with its parameters elided.}
\pair{machine-learning/first/7}{Apply the model to a point to get a class.}
\pair{machine-learning/first/8}{With \mcode{"Probabilities"} it returns one rule
per class.}
\pair{machine-learning/first/9}{\B{Map} applies it to several points. The cloud
centres come back as their own classes.}
\pair{machine-learning/first/10}{The model answers questions about itself.}
\end{codepairs}

The default method is \emph{nearest neighbours} with one neighbour, the subject of
the next section. The classes are kept in the order they first appear in the
training data, here \mcode{"A"}, \mcode{"B"}, \mcode{"C"}. That order is also the
order of the probability rules, and it breaks ties: when two classes are equally
likely, the one that appeared first wins. None of the methods sorts the classes.

A class does not have to be a string. It can be any expression: a symbol, a number,
or even a list. Classes are compared structurally, the way \B{SameQ} compares
them, so \mcode{"a"} and \mcode{a} are two different classes.

\begin{codepairs}
\pair{machine-learning/forms/1}{Three classes: the symbol \mcode{x}, the list
\mcode{\{1, 2\}}, and the string \mcode{"y"}.}
\pair{machine-learning/forms/2}{The class vocabulary, in order of first
appearance.}
\pair{machine-learning/forms/3}{A list is returned as a class like any other.}
\pair{machine-learning/forms/4}{\B{FullForm} shows the whole model: the method, a
payload, the feature count, and \(k\).}
\end{codepairs}

The payload is the class vocabulary followed by one row per training example. Each
row holds the features and then the \emph{index} of the class, starting from
zero. So the numeric part of the model is purely numeric, and the vocabulary is the
one place a class is named.\index{classifier!label vocabulary}

\calloutnote{Under the hood}{A fitted model is not a new kind of object. It is an
ordinary expression whose head is \mcode{ClassifierFunction},
\B{PredictorFunction}, \B{DimensionReducerFunction}, or
\mcode{LearnedDistribution}, and whose four arguments are a method name, a payload
of plain lists, and two integers. Applying one uses the evaluator's existing
dispatch for compound heads (\texttt{ml\_model\_apply} in
\texttt{src/ml/predict.c}, called from \texttt{src/eval.c}), which is the same
mechanism that applies a pure function or an association. The design notes in
\texttt{src/ml/predict.h} record why: a new node type would need new cases in
copying, freeing, hashing, comparison, and printing, to store numbers that lists
already hold. The printer elides the payload because a nearest-neighbour model
carries its entire training set.}

\subsection{What Classify does not accept}
\label{ssec:ml-limits}

Mathematica's \mcode{Classify} accepts several data layouts, features of mixed
type, and a batch of points in one call. Mathilda's accepts one layout and numeric
features only. An unsupported call returns unevaluated rather than guessing:

\begin{codepairs}
\pair{machine-learning/forms/5}{The \(\text{list} \to \text{list}\) layout is not
accepted. Use \mcode{Thread} to make rules.}
\pair{machine-learning/forms/6}{A string feature is refused. Features must be
real numbers.}
\pair{machine-learning/forms/7}{A matrix of points is not a batch. Use \mcode{Map}
instead.}
\pair{machine-learning/forms/8}{An option that means nothing to the method is
refused, not ignored.}
\pair{machine-learning/forms/9}{A method that is not implemented is refused.}
\end{codepairs}

\calloutnote{Pitfall}{Because every feature must be a number, a categorical
feature such as a colour has to be encoded by hand before training. The usual
choice is \emph{one-hot} encoding: one 0/1 feature per category. Coding the
categories as \(1, 2, 3\) instead imposes an order and a spacing that the
categories do not have, and a distance-based method will take that spacing
seriously.}

\section{Nearest neighbours}
\label{sec:ml-knn}

The \(k\)-nearest-neighbour classifier\index{nearest neighbours!classification}
has no training step at all. The training set is the model. To classify a query
point \(x\), it computes the squared Euclidean distance
\[
  d(x, x_i) \;=\; \sum_{j=1}^{p} (x_j - x_{ij})^2
\]
to every training point \(x_i\), keeps the \(k\) smallest, and lets those \(k\)
points vote. The probability reported for class \(c\) is its share of the vote,
\(\#\{\,i \in N_k(x) : y_i = c\,\}/k\), where \(N_k(x)\) is the set of the \(k\)
nearest training points. Squaring the distance does not change which points are
nearest, and it saves a square root per comparison.

With \(k=1\), the default, each training point claims the region of the plane
closer to it than to any other training point. These regions are the cells of the
\emph{Voronoi diagram}\index{Voronoi diagram} of the training set, and the decision
boundary is made of cell edges. The plot below colours a fine grid of query points
by the class the classifier assigns and draws the training points on top.

\plotcell[0.6\linewidth]{machine-learning/knn-1}

The figure sessions label the classes \(1, 2, 3\) so that a class can index a list
of colours. The helper \mcode{picture} classifies about \(4600\) grid points and
draws each one as a small rectangle. The same helper draws every decision-region
figure in this chapter, so the figures can be compared directly.

With one neighbour every training point is classified correctly, including the
orange point that sits inside the blue cloud near \((0.8, 1.1)\). That point claims
a small island of its own. This is \emph{overfitting}\index{overfitting}: the
boundary follows the noise in the training set. More neighbours smooth it out.
With fifteen neighbours the islands disappear and the boundaries straighten:

\plotcell[0.6\linewidth]{machine-learning/knn-15}

The number of neighbours is a sub-option of the method:

\begin{codepairs}
\pair{machine-learning/knn/4}{Five neighbours vote.}
\pair{machine-learning/knn/5}{At \((1.5, 2)\), three of the five nearest points
are from \mcode{"B"}.}
\end{codepairs}

Section~\ref{sec:ml-eval} measures how the choice of \(k\) affects accuracy on
new data.

\subsection{Scale matters}
\label{ssec:ml-scale}

A distance adds up the squared differences of all the features, so a feature
measured on a large scale dominates the others. Suppose the first coordinate were
recorded in millimetres and the second in metres. Multiplying the first coordinate
by \(1000\) leaves the problem unchanged, but the classifier now ignores the second
coordinate almost entirely:

\begin{codepairs}
\pair{machine-learning/knn/6}{The first feature, rescaled by \(1000\).}
\pair{machine-learning/knn/7}{The centre of cloud \mcode{"C"} is now classified as
\mcode{"A"}, because only the first coordinate counts.}
\pair{machine-learning/knn/8}{The training means and standard deviations of each
column.}
\pair{machine-learning/knn/9}{\B{Standardize} gives each column mean \(0\) and
standard deviation \(1\).}
\pair{machine-learning/knn/10}{A new point is scaled with the \emph{training}
statistics before it is classified. The answer is \mcode{"C"} again.}
\end{codepairs}

\mcode{Standardize} computes \((x - \bar{x})/s\) for each column, with the sample
standard deviation \(s\) (divisor \(n-1\)), so it agrees with \B{Mean} and
\B{StandardDeviation}. A query point must be scaled with the training means and
deviations, not its own. \mcode{Standardize} returns only the scaled data, so the
session keeps \mcode{mu} and \mcode{sd} to apply to new points.

\calloutnote{Under the hood}{The classifier finds the nearest points by a linear
scan of the payload, keeping the \(k\) best in a small sorted array
(\texttt{ml\_classifier\_apply} in \texttt{src/ml/classify.c}). Ties in distance
keep the earlier training row, and ties in the vote go to the class that appeared
first, so the answer never depends on scan order. There is no spatial index, so
each query costs \(O(np)\) for \(n\) training points with \(p\) features. Training
costs almost nothing, and prediction costs a pass over the data.}

\section{Naive Bayes}
\label{sec:ml-bayes}

A \emph{probabilistic} classifier models how each class generates its points and
then inverts the model with Bayes' rule.\index{Bayes' rule} If \(\pi_c\) is the
prior probability of class \(c\) and \(f_c(x)\) is the density of its points, the
posterior probability of class \(c\) at \(x\) is
\[
  P(c \mid x) \;=\; \frac{\pi_c\, f_c(x)}{\sum_{c'} \pi_{c'}\, f_{c'}(x)} .
\]
\emph{Gaussian naive Bayes}\index{naive Bayes} takes each \(f_c\) to be a product
of one-dimensional normal densities, one per feature:
\[
  f_c(x) \;=\; \prod_{j=1}^{p} \frac{1}{\sqrt{2\pi\sigma_{cj}^2}}
    \exp\!\left(-\frac{(x_j - \mu_{cj})^2}{2\sigma_{cj}^2}\right).
\]
The product is the ``naive'' assumption: within a class, the features are treated
as independent. Equivalently, each class has a diagonal covariance matrix.
Training estimates \(\pi_c\) as the fraction of training points in class \(c\),
and \(\mu_{cj}\) and \(\sigma^2_{cj}\) as the mean and variance of feature \(j\)
within class \(c\). That is one pass over the data.

A one-feature example shows the whole model. Eight points on a line, with
\mcode{"lo"} mostly on the left and \mcode{"hi"} mostly on the right:

\begin{codepairs}
\pair{machine-learning/bayes/1}{The point at \(4\) is \mcode{"hi"} and the point
at \(5\) is \mcode{"lo"}, so the classes overlap.}
\pair{machine-learning/bayes/2}{A naive Bayes fit.}
\pair{machine-learning/bayes/3}{The payload: vocabulary, priors, and then a mean
row and a variance row for each class.}
\pair{machine-learning/bayes/4}{The posterior at \(3.9\).}
\pair{machine-learning/bayes/5}{The same posterior from Bayes' rule and \B{PDF},
written out by hand.}
\pair{machine-learning/bayes/6}{The two variances are equal, so the boundary is
the midpoint \(4.5\). A tie goes to the first class.}
\end{codepairs}

The class \mcode{"lo"} holds \(\{1,2,3,5\}\), with mean \(2.75\) and variance
\(2.1875\). The class \mcode{"hi"} holds \(\{4,6,7,8\}\), with mean \(6.25\) and
the same variance. The variance uses the divisor \(n_c\) (the maximum-likelihood
estimate), not \(n_c - 1\). The hand computation in line 5 uses \mcode{PDF} of a
\B{NormalDistribution}, which is separate code from the classifier, and the two
agree to every printed digit.

When the class variances differ, the boundary is no longer a single point. Taking
logarithms, the classifier compares
\[
  \log \pi_c \;-\; \frac{1}{2}\sum_{j} \left( \log (2\pi\sigma_{cj}^2)
    + \frac{(x_j-\mu_{cj})^2}{\sigma_{cj}^2} \right)
\]
across classes, and the difference between two classes is a quadratic function of
\(x\). In one dimension a quadratic can have two roots, so the class with the
larger variance can win on \emph{both} sides:

\begin{codepairs}
\pair{machine-learning/bayes/7}{Now \mcode{"lo"} is tight around \(2\) and
\mcode{"hi"} is spread from \(5\) to \(13\).}
\pair{machine-learning/bayes/8}{Far to the left, at \(-10\), the wide class wins
again.}
\end{codepairs}

This is correct for the model. Far from both means, the density with the heavier
spread is the larger one. In two dimensions the boundaries are conic sections,
which the running example shows clearly:

\plotcell[0.6\linewidth]{machine-learning/bayes-regions}

\calloutnote{Under the hood}{The classifier works with log densities and
converts to probabilities once, with a softmax: it subtracts the largest log
posterior before exponentiating. A product of many small densities underflows to
zero in floating point for a point several standard deviations out, and then every
class would look equally likely at exactly the points where the answer is
clearest.}

\calloutnote{Pitfall}{If a feature is constant within one class, its variance
there is zero and its density is infinite at that value. Mathilda puts a floor
under every variance: \(10^{-6}\) times the variance of that feature over the
whole training set (\texttt{src/ml/classify.c}). The floor scales with the
feature, so it behaves the same whether the feature is in millimetres or
kilometres. It is still tiny, so a feature that never varies in one class acts
almost as a veto. The spam example in \S\ref{sec:ml-featured} runs into exactly
this.}

\section{Logistic regression}
\label{sec:ml-logistic}

Naive Bayes models the points and derives the class probabilities.
\emph{Logistic regression}\index{logistic regression} models the probabilities
directly. For two classes it sets
\[
  P(y = 1 \mid x) \;=\; \sigma(\beta_0 + \beta^{\mathsf T} x),
  \qquad \sigma(t) \;=\; \frac{1}{1 + e^{-t}},
\]
and chooses \(\beta\) to maximise the log-likelihood
\[
  \ell(\beta) \;=\; \sum_{i=1}^{n} \Bigl( y_i \log p_i + (1-y_i)\log(1-p_i) \Bigr),
  \qquad p_i = \sigma(\beta_0 + \beta^{\mathsf T} x_i).
\]
The decision boundary \(P = 1/2\) is the hyperplane
\(\beta_0 + \beta^{\mathsf T} x = 0\). Here \(y = 1\) means the second class in
the vocabulary.

\(\ell\) is concave, so Newton's method finds its maximum. Each Newton step turns
out to be a weighted least-squares problem, which is why the method is called
\emph{iteratively reweighted least squares} (IRLS).\index{IRLS} With the design
matrix \(A = [\,1 \mid X\,]\), the weights \(w_i = p_i(1 - p_i)\), and the working
response \(z_i = \eta_i + (y_i - p_i)/w_i\) where \(\eta_i\) is the current linear
predictor, the next estimate solves
\[
  (A^{\mathsf T} W A)\,\beta \;=\; A^{\mathsf T} W z .
\]
The one-feature example from the previous section:

\begin{codepairs}
\pair{machine-learning/logistic/1}{The same eight points.}
\pair{machine-learning/logistic/2}{A logistic fit.}
\pair{machine-learning/logistic/3}{The coefficients \(\beta_0, \beta_1\), read from
the payload.}
\pair{machine-learning/logistic/4}{The boundary \(-\beta_0/\beta_1\) is exactly
\(4.5\).}
\pair{machine-learning/logistic/5}{And the probability there is exactly one
half.}
\pair{machine-learning/logistic/6}{The logistic function by hand, evaluated at
\(x = 6\).}
\pair{machine-learning/logistic/7}{This value is the probability the model assigns to \mcode{"hi"}.}
\end{codepairs}

The data are symmetric about \(4.5\), so the fitted boundary lands there exactly.
Line 5 is the check that ties the fit and the application together: any mistake
in how the coefficients are stored or read back would move the probability away
from \(1/2\). The fitted curve is the familiar S shape:

\plotcell[0.6\linewidth]{machine-learning/logistic-curve}

\calloutnote{Theory}{If the two classes can be separated by a hyperplane, the
likelihood has no maximum. Scaling \(\beta\) up sends every \(p_i\) towards \(0\)
or \(1\) and \(\ell\) towards \(0\), so plain Newton iteration runs off to
infinity. Mathilda adds a small ridge, \(10^{-6}\), to the diagonal of
\(A^{\mathsf T} W A\) for the non-intercept coefficients. This is the Newton step
for a penalised likelihood, which is strictly concave, so the fit is finite and
unique. The intercept is not penalised, because shrinking it would bias the base
rate. The iteration also stops after 100 steps, and each weight is clamped at
\(10^{-10}\) so that the working response never divides by zero.}

Separable data therefore give a finite answer, with the boundary in the middle of
the gap:

\begin{codepairs}
\pair{machine-learning/logistic/8}{Two points each side of a wide gap.}
\pair{machine-learning/logistic/9}{Large but finite coefficients.}
\pair{machine-learning/logistic/10}{The probability is one half at \(5\), the
middle of the gap.}
\end{codepairs}

\subsection{More than two classes}
\label{ssec:ml-ovr}

For \(K > 2\) classes Mathilda fits \emph{one-vs-rest}:\index{one-vs-rest} \(K\)
separate binary models, the \(k\)th treating class \(k\) as positive and all the
others as negative. The predicted class is the one whose model gives the largest
probability. On the running example the boundaries are straight line segments:

\plotcell[0.6\linewidth]{machine-learning/logistic-regions}

One-vs-rest has a weakness that a small example exposes. Put three classes in a
row on a line:

\begin{codepairs}
\pair{machine-learning/logistic/11}{Three classes in a row.}
\pair{machine-learning/logistic/12}{Three coefficient vectors. The middle one is
flat: slope \(0\), intercept \(\log(1/2)\).}
\pair{machine-learning/logistic/13}{In the middle, \mcode{"b"} still wins.}
\end{codepairs}

No single threshold separates \mcode{"b"} from \mcode{"a"} and \mcode{"c"}
together, so the best one-feature logistic model for ``\mcode{"b"} against the
rest'' is a constant. Its intercept \(\log(1/2)\) gives \(\sigma = 1/3\), the
fraction of \mcode{"b"} in the data. The class still comes out right, because the
models for \mcode{"a"} and \mcode{"c"} are both close to zero in the middle.

\calloutnote{Pitfall}{The \(K\) probabilities of a one-vs-rest model come from
\(K\) independent fits, and nothing makes them add up to one. Mathilda divides
them by their sum before reporting them. Division by a common factor does not
change which is largest, so the predicted class is exactly what the \(K\) models
say, but the reported probabilities are a convention and not calibrated
likelihoods. The source (\texttt{src/ml/classify.c}) gives the reason for
one-vs-rest over a softmax model: a softmax's coefficients are defined only up to
an additive constant, so they could not be checked by a test.}

\section{Decision trees and random forests}
\label{sec:ml-trees}

A \emph{decision tree}\index{decision tree} classifies by asking a sequence of
yes-or-no questions of the form ``is feature \(j\) at most \(t\)?''. Each answer
sends the point to one of two subtrees, and a leaf gives the answer. Mathilda grows
trees by CART, the \emph{classification and regression trees} algorithm of
Breiman, Friedman, Olshen, and Stone. At each node it chooses the split that most
reduces the \emph{Gini impurity}\index{Gini impurity}
\[
  G \;=\; 1 - \sum_{c} q_c^2 ,
\]
where \(q_c\) is the fraction of the node's points in class \(c\). A pure node has
\(G = 0\). A split of \(m\) points into children of sizes \(m_L\) and \(m_R\) is
scored by the weighted impurity \(m_L G_L + m_R G_R\), and the best split has the
smallest.

\begin{codepairs}
\pair{machine-learning/trees/1}{The eight points again.}
\pair{machine-learning/trees/2}{A CART tree.}
\pair{machine-learning/trees/3}{Seven nodes. Each row is \{feature, threshold,
left, right\}. A feature of \(-1\) marks a leaf.}
\pair{machine-learning/trees/4}{The class counts at every node, in vocabulary
order \(\{\)\mcode{"lo"}, \mcode{"hi"}\(\}\).}
\pair{machine-learning/trees/5}{The tree carves out the single point \(5\) as
\mcode{"lo"}.}
\end{codepairs}

Node and feature numbers start at zero. Read the table from the top. Node 0 asks
whether \(x \le 3.5\). If so, the point goes to node 1, a leaf holding the three
\mcode{"lo"} points. If not, node 2 asks whether \(x \le 5.5\), and so on down to
node 4 (the point \(4\), \mcode{"hi"}) and node 5 (the point \(5\),
\mcode{"lo"}). Every threshold lies halfway between two neighbouring training
values: \(3.5\), \(4.5\), \(5.5\). A threshold equal to a training value would
make the side a new point lands on depend on floating-point equality, while a
midpoint puts the boundary in the gap where no training point lies.

Mathilda grows the tree until every leaf is pure or cannot be split (the depth
limit is 32). So a tree reproduces its training labels exactly, like one nearest
neighbour, and it overfits in the same way. On the running example, a 19-node
tree gets all ninety training points right:

\begin{codepairs}
\pair{machine-learning/trees/6}{The running example.}
\pair{machine-learning/trees/7}{}
\pair{machine-learning/trees/8}{}
\pair{machine-learning/trees/9}{A 19-node tree.}
\pair{machine-learning/trees/10}{All ninety training points are classified
correctly.}
\end{codepairs}

Every boundary a tree draws is parallel to an axis, because every question
involves one feature. The regions are unions of rectangles, and the narrow strips
are the price of isolating individual points:

\plotcell[0.6\linewidth]{machine-learning/tree-regions}

\calloutnote{Under the hood}{For each node and feature, the splitter
(\texttt{src/ml/tree.c}) sorts the node's points by that feature. It then walks
the sorted order, moving one point at a time from the right child to the left and
updating both class histograms. The impurity is evaluated only where the feature
value changes. That costs \(O(m \log m)\) per feature per node for the sort and
\(O(m)\) for the sweep, where recomputing the histograms for every candidate
threshold would cost \(O(m^2)\). Ties in impurity go to the lower feature index and
then the lower threshold, and the sort breaks equal values by point index, so the
same data always give the same tree. The source chooses Gini over entropy because
it needs no logarithm and has no special case at \(q_c = 0\).}

\subsection{Random forests}
\label{ssec:ml-forest}

A single deep tree has low bias and high variance: a small change in the data can
change the tree a lot. A \emph{random forest}\index{random forest} averages many
trees to reduce the variance. Mathilda grows fifty trees. Each tree sees a
\emph{bootstrap} resample of the training set, \(n\) rows drawn with replacement,
and each node considers only a random subset of \(\operatorname{round}(\sqrt{p})\)
of the \(p\) features. The forest reports the average of the trees' class
distributions.

\begin{codepairs}
\pair{machine-learning/trees/11}{Two forests grown from the same seed.}
\pair{machine-learning/trees/12}{They are identical. A forest holds fifty
trees.}
\pair{machine-learning/trees/13}{The trees differ in size, because each saw
different data.}
\pair{machine-learning/trees/14}{The trees disagree near a boundary, so the
probabilities are mixed.}
\pair{machine-learning/trees/15}{Another seed gives another forest and slightly
different probabilities.}
\end{codepairs}

A single tree on this data gives probabilities of \(0\) or \(1\) almost everywhere.
The forest's averages carry real information about how confident the ensemble is.
Its decision regions are smoother than one tree's, but they keep the axis-aligned
character:

\plotcell[0.6\linewidth]{machine-learning/forest-regions}

\calloutnote{Theory}{Averaging \(B\) predictors, each with variance \(\sigma^2\)
and pairwise correlation \(\rho\), gives variance
\(\rho\sigma^2 + (1-\rho)\sigma^2/B\). More trees shrink only the second term.
The first term needs the trees to be less correlated, and that is the job of the
feature sampling. If every tree could choose among all features, most would split
first on the same strong feature and look alike. Restricting each node to a
random subset forces variety. With two features, as here, each node considers
just one.}

\calloutnote{Under the hood}{Both the bootstrap and the feature sampling draw from
\texttt{random\_uniform\_01}, the generator behind \mcode{RandomReal}. This is why
\mcode{SeedRandom} makes a forest reproducible, as line 12 checks. Each tree's leaf
counts are normalised to a distribution before the fifty are averaged. Pooling the
raw counts instead would give a tree with large leaves more weight than a tree
with small ones. The forest size is fixed at fifty. Mathematica tunes it on
held-out data, which would mean setting aside part of your data without asking.}

\section{Measuring a classifier}
\label{sec:ml-eval}

Accuracy on the training data says little. One nearest neighbour and a full tree
score \(100\%\) on any training set without duplicated contradictory points. The
honest measure is accuracy on data the model has not seen. The session below draws
a fresh \emph{test set} of 300 points from the same three clouds and measures each
method on it:\index{test set}

\begin{codepairs}
\pair{machine-learning/evaluate/1}{Rebuild the running example's training set.}
\pair{machine-learning/evaluate/2}{}
\pair{machine-learning/evaluate/3}{}
\pair{machine-learning/evaluate/4}{Then draw 300 fresh points from the same
clouds, with a different seed.}
\pair{machine-learning/evaluate/5}{Accuracy: the fraction of examples classified
correctly.}
\pair{machine-learning/evaluate/6}{Perfect on the training data, \(94.7\%\) on the
test set.}
\pair{machine-learning/evaluate/7}{Pairs of \{true class, predicted class\}.}
\pair{machine-learning/evaluate/8}{The confusion matrix: rows are true classes,
columns are predictions.}
\end{codepairs}

The \emph{confusion matrix}\index{confusion matrix} shows where the errors are.
Cloud \mcode{"A"} is almost never confused with anything. Most errors are
\mcode{"C"} points called \mcode{"B"}, where the two clouds overlap. Mathilda has
no \mcode{ClassifierMeasurements} function, but \B{Count} over the pairs gives the
matrix, and any other measure (precision, recall, per-class accuracy) is one line
more.

Now compare the methods, on the training data and on the test set:

\begin{codepairs}
\pair{machine-learning/evaluate/9}{All five methods.}
\pair{machine-learning/evaluate/10}{\{method, training accuracy, test accuracy\}
for each.}
\end{codepairs}

One nearest neighbour, perfect in training, is the worst on the test set.
Naive Bayes, which makes errors in training, is the best. The data
really are Gaussian clouds, so naive Bayes has the right model. On data that are
not Gaussian, the order would change.

\subsection{Choosing \texorpdfstring{\(k\)}{k}}
\label{ssec:ml-choosek}

The number of neighbours trades variance for bias. Too few, and the classifier
follows noise. Too many, and it averages over points from other classes. The plot
shows test accuracy for \(k = 1, \dots, 60\):

\plotcell[0.6\linewidth]{machine-learning/k-sweep}

Accuracy rises from \(94.7\%\) at \(k=1\) to about \(97\%\) by \(k=5\),
stays there up to \(k \approx 57\), and then starts to fall as each vote reaches
into neighbouring classes. The classes are well separated, so the plateau is
long; on harder data it is shorter. The curve is jagged because a test set of 300
points changes accuracy in steps of \(1/300\).

In practice there is no separate test set, and \(k\) has to be chosen from the
training data. \emph{Cross-validation}\index{cross-validation} does this. Split the
training data into \(m\) parts, train on \(m-1\) of them, test on the remaining
one, and average over the \(m\) choices. Mathilda has no built-in cross-validation,
but it takes two lines:

\begin{codepairs}
\pair{machine-learning/evaluate/11}{Five folds of eighteen, from a random
permutation.}
\pair{machine-learning/evaluate/12}{Train on four folds, test on the fifth, and
average.}
\pair{machine-learning/evaluate/13}{Cross-validated accuracy of the four
deterministic methods.}
\end{codepairs}

The ranking agrees with the test set: the two methods that assume smooth
boundaries (naive Bayes and logistic regression) beat the two that can fit
anything (one neighbour and a full tree). With only ninety points, the numbers
move by a few percent with the random split.

\subsection{When the boundary is not straight}
\label{ssec:ml-rings}

Every method so far has done well on the running example because its classes are
roughly separable by straight lines. The next data set is not. Class
\mcode{"in"} is a noisy ring of radius \(1\), and class \mcode{"out"} a ring of
radius \(2.5\) around it. No straight line separates them:

\begin{codepairs}
\pair{machine-learning/rings/1}{Noisy points on a circle.}
\pair{machine-learning/rings/2}{Two hundred points make the training set.}
\pair{machine-learning/rings/3}{Six hundred more make the test set.}
\pair{machine-learning/rings/4}{}
\pair{machine-learning/rings/5}{Logistic regression scores \(50\%\), no better
than a coin.}
\end{codepairs}

Nearest neighbours and the tree adapt to any shape. Naive Bayes also works well,
and the reason is instructive. Both classes are centred at the origin, but the
outer class has a much larger variance in each coordinate. The naive Bayes
boundary is a conic section, and here the conic is a closed curve around the
origin:

\plotcell[0.55\linewidth]{machine-learning/rings-bayes}

Logistic regression can only draw a straight line, and it cannot do better than
chance. It does not need a new method, though. It needs a better feature. The
squared radius \(r^2 = x_1^2 + x_2^2\) separates the rings with a single threshold,
so add it as a third feature:\index{feature engineering}

\begin{codepairs}
\pair{machine-learning/rings/6}{Append \(x \cdot x\) to each feature vector.}
\pair{machine-learning/rings/7}{Logistic regression on three features.}
\pair{machine-learning/rings/8}{\(99.8\%\) on the test set.}
\pair{machine-learning/rings/9}{The weight on \(r^2\) dominates, and the weights on
\(x_1, x_2\) are small.}
\pair{machine-learning/rings/10}{The boundary is the circle
\(r = \sqrt{-\beta_0/\beta_3}\), close to \(1.75\), halfway between the rings.}
\end{codepairs}

A linear model in the new features is a quadratic model in the old ones. This is
the idea behind kernel methods and behind much of practical machine learning:
choosing features well often matters more than choosing the method.

\section{Regression with Predict}
\label{sec:ml-predict}

\mcode{Predict} learns a numerical response instead of a class. Its default method is
\emph{least squares}\index{least squares}: it fits
\(y \approx \beta_0 + \beta^{\mathsf T} x\) by minimising
\(\sum_i (y_i - \beta_0 - \beta^{\mathsf T} x_i)^2\). With the design matrix
\(A = [\,1 \mid X\,]\), the minimiser solves the \emph{normal equations}
\[
  A^{\mathsf T} A\, \beta \;=\; A^{\mathsf T} y .
\]
Forty noisy points on the line \(y = 2.5 + 0.8x\):

\begin{codepairs}
\pair{machine-learning/predict/1}{Forty points: true intercept \(2.5\), slope
\(0.8\), and unit noise.}
\pair{machine-learning/predict/2}{A \mcode{PredictorFunction}.}
\pair{machine-learning/predict/3}{The coefficients \(\beta_0, \beta_1\), and a
prediction at \(x = 5\).}
\pair{machine-learning/predict/4}{\B{Fit} computes the same line by its own
route.}
\pair{machine-learning/predict/5}{\B{LeastSquares} on the design matrix gives the
same coefficients.}
\pair{machine-learning/predict/6}{A matrix whose last column is the response is
accepted too.}
\end{codepairs}

Three separate routes give the same coefficients to six digits. The estimates
\(2.40\) and \(0.82\) differ from the true \(2.5\) and \(0.8\) because of the
noise, which is what a finite sample does. With several features the
coefficients come back in feature order, after the intercept. Noise-free data
from \(y = 1 + 2x_1 + 3x_2\) recover the coefficients exactly:

\begin{codepairs}
\pair{machine-learning/predict/7}{Intercept \(1\), then \(2\) and \(3\).}
\pair{machine-learning/predict/8}{The second feature is twice the first, so the fit
is not unique, and \mcode{Predict} declines.}
\end{codepairs}

In line 8 the second feature is exactly twice the first, so infinitely many
coefficient vectors fit equally well. Mathilda returns the call unevaluated
instead of picking one.

\calloutnote{Under the hood}{\texttt{ml\_ols} in \texttt{src/ml/predict.c} forms
\(A^{\mathsf T} A\) and \(A^{\mathsf T} y\) in one pass, without building \(A\),
and solves by Gaussian elimination with partial pivoting. The comment there
explains why it does not use LAPACK's QR-based \texttt{dgels}, which is
numerically better because it avoids squaring the condition number. The system
has only \(p+1\) rows for \(p\) features, and a zero pivot in elimination reports
collinear features directly, which is what makes the refusal in line 8 possible.
With many nearly collinear features the normal equations lose accuracy, and
\mcode{LeastSquares} is the better tool.}

\subsection{Nearest-neighbour regression}
\label{ssec:ml-knnreg}

\mcode{Predict} with \mcode{Method -> "NearestNeighbors"} predicts the mean response
of the \(k\) nearest training points. The default is \(k = 3\). A regression
averages, so a few neighbours smooth out noise, while a classifier votes, and its
default of one neighbour reproduces the training labels exactly.

\begin{codepairs}
\pair{machine-learning/predict/9}{A nearest-neighbour predictor.}
\pair{machine-learning/predict/10}{The prediction at \(5\), and \(k\).}
\pair{machine-learning/predict/11}{The same number by hand: sort by distance to
\(5\), take three, average.}
\pair{machine-learning/predict/12}{Other methods, such as random forests, are not
implemented.}
\end{codepairs}

The two predictors side by side: the least-squares line in blue and the
three-neighbour predictor in orange. The nearest-neighbour curve is a step
function, constant wherever the set of three nearest points does not change. It
follows the local wiggles of the data, noise included:

\plotcell[0.6\linewidth]{machine-learning/regression}

\calloutnote{Pitfall}{\B{LinearModelFit} exists, but only as another name for the
least-squares \mcode{Predict}. It returns a \mcode{PredictorFunction} with the same
coefficients. Mathematica's \mcode{LinearModelFit} returns a \mcode{FittedModel}
with regression diagnostics such as \(R^2\), standard errors, and an ANOVA table.
None of those are implemented, and the source says so rather than approximating
them.}

\section{Finding clusters}
\label{sec:ml-clusters}

Clustering\index{clustering} is classification without labels: the task is to
find groups in the data. \mcode{FindClusters} takes a list of numbers, of vectors, of
colours, or of strings, and returns the groups as a list of lists. It keeps the
elements in their input order, and orders the clusters by first appearance. On a
line the answer is easy to follow:

\begin{codepairs}
\pair{machine-learning/clusters/1}{Three groups, found automatically.}
\pair{machine-learning/clusters/2}{An exact count, and an upper bound with
\mcode{UpTo}.}
\end{codepairs}

With exactly two clusters the two closest groups merge. With \mcode{UpTo[4]},
\mcode{FindClusters} returns the three groups the data suggest, not four.

The default method is \emph{single linkage}\index{single linkage}: repeatedly
merge the two clusters whose closest members are nearest. The result is the same
as building the \emph{minimum spanning tree} of the points and deleting its
longest edges. With no count given, Mathilda cuts every tree edge longer than
three times the median edge. Single linkage finds clusters of any shape, as long
as a gap separates them. It fails when a chain of points bridges two clusters, or
when an outlier sits far from everything:

\begin{codepairs}
\pair{machine-learning/clusters/3}{The running example, without labels.}
\pair{machine-learning/clusters/4}{}
\pair{machine-learning/clusters/5}{Single linkage finds one big cluster and an
outlier.}
\pair{machine-learning/clusters/6}{Asked for three, it still splits off
outliers.}
\pair{machine-learning/clusters/7}{\(k\)-means finds three clusters of about thirty
each.}
\pair{machine-learning/clusters/8}{Their means are close to the true centres
\((0,0)\), \((3,3)\), \((0,4)\).}
\end{codepairs}

The three clouds overlap, so there is no gap between them, and single linkage
cannot separate them. \emph{\(k\)-means}\index{k-means} solves a different
problem. It looks for \(k\) centres \(m_1, \dots, m_k\) that minimise the total
squared distance from each point to its nearest centre,
\[
  \sum_{i=1}^{n} \min_{j} \lVert x_i - m_j \rVert^2 .
\]
Lloyd's algorithm alternates two steps. It assigns each point to its nearest
centre, then moves each centre to the mean of its points. Neither step can
increase the objective, so the iteration converges, though possibly to a local
minimum. Here are the three clusters it finds, coloured by cluster:

\plotcell[0.6\linewidth]{machine-learning/kmeans-blobs}

\calloutnote{Under the hood}{Mathilda's \(k\)-means starts deterministically, with
\emph{farthest-first} seeding: the first centre is the data point nearest the
overall mean, and each further centre is the point farthest from the centres
already chosen. So the result does not depend on the order of the input, and no
random seed is needed. A centre that loses all its points is moved to the point
farthest from its own centre. \mcode{"KMeans"} needs a count, either exactly \(k\)
or \mcode{UpTo[k]}, and returns unevaluated without one (line 9 below). The
details are in \texttt{src/list/find\_clusters.c} and its reference page.}

The two methods fail in opposite ways. On two concentric rings, \(k\)-means cuts
the plane with a straight line and splits each ring in half:

\plotcell[0.5\linewidth]{machine-learning/rings-kmeans}

Single linkage follows each ring around, because consecutive points on a ring are
closer to each other than the rings are:

\plotcell[0.5\linewidth]{machine-learning/rings-linkage}

\begin{table}[ht]
\centering
\small
\begin{tabular}{@{}llccc@{}}
\toprule
Method & Idea & automatic & \mcode{UpTo[n]} & \(n\) \\
\midrule
\mcode{"Agglomerate"} (default) & single linkage & yes & yes & yes \\
\mcode{"SpanningTree"} & minimum spanning tree & yes & yes & yes \\
\mcode{"KMeans"} & Lloyd's algorithm & no & yes & yes \\
\mcode{"KMedoids"} & centres are data points & no & yes & yes \\
\mcode{"Spectral"} & graph Laplacian & yes & yes & no \\
\mcode{"DBSCAN"} & density-connected regions & yes & no & no \\
\mcode{"GaussianMixture"} & EM, count by BIC & yes & no & no \\
\mcode{"JarvisPatrick"} & shared nearest neighbours & yes & no & no \\
\mcode{"MeanShift"} & modes of a density & yes & no & no \\
\mcode{"NeighborhoodContraction"} & moving points to local means & yes & no & no \\
\bottomrule
\end{tabular}
\caption{The ten \mcode{FindClusters} methods, and which ways of giving the count
each accepts.}
\label{tab:ml-cluster-methods}
\end{table}

Table~\ref{tab:ml-cluster-methods} lists all ten methods. All of them work on
vectors of any dimension. Methods with a neighbourhood radius (\mcode{"DBSCAN"},
\mcode{"MeanShift"}, \mcode{"NeighborhoodContraction"}) take it as a sub-option,
and their results depend strongly on it. Mathematica chooses a distance function
and preprocesses the data by its own unpublished rules, so its cluster
assignments will often differ from Mathilda's.

\subsection{Classifying new points after clustering}
\label{ssec:ml-clusterclassify}

Mathematica's \mcode{ClusterClassify} returns a function that assigns new points
to the clusters it found. Mathilda has no such function, but a clustering is a
labelling, and \mcode{Classify} can learn it:

\begin{codepairs}
\pair{machine-learning/clusters/9}{\(k\)-means needs a count.}
\pair{machine-learning/clusters/10}{Label each point with the number of its
cluster.}
\pair{machine-learning/clusters/11}{Train a classifier on those labels.}
\pair{machine-learning/clusters/12}{New points are assigned to clusters.}
\end{codepairs}

A nearest-neighbour classifier trained this way assigns a new point to the
cluster of its nearest clustered point. For \(k\)-means you could instead compute
the nearest cluster mean, which gives the Voronoi cells of the centres.

\section{Gaussian mixtures}
\label{sec:ml-mixtures}

\(k\)-means measures distance the same way in every direction, so it assumes the
clusters are round and of similar size. The next two clouds are long, thin, and
tilted towards each other. Each has 150 points:

\begin{codepairs}
\pair{machine-learning/mixture/1}{Start with standard normal pairs.}
\pair{machine-learning/mixture/2}{They are stretched and sheared into two tilted
clouds.}
\pair{machine-learning/mixture/3}{\(k\)-means with two clusters, and a Gaussian
mixture.}
\pair{machine-learning/mixture/4}{Cluster sizes.}
\pair{machine-learning/mixture/5}{How many of the first cloud's 150 points land in
the first cluster of each.}
\end{codepairs}

\(k\)-means splits the data along the wrong line. Only 23 points of its first
cluster come from the first cloud:

\plotcell[0.5\linewidth]{machine-learning/gmm-kmeans}

A \emph{Gaussian mixture}\index{Gaussian mixture} models the data as coming from
\(K\) normal distributions, each with its own mean \(\mu_j\), full covariance
matrix \(\Sigma_j\), and weight \(w_j\):
\[
  f(x) \;=\; \sum_{j=1}^{K} w_j\, \mathcal{N}(x;\, \mu_j, \Sigma_j).
\]
The mixture puts all but two of the first cloud's points in one cluster:

\plotcell[0.5\linewidth]{machine-learning/gmm-clusters}

The parameters are fitted by the \emph{expectation--maximisation}
(EM)\index{expectation--maximisation} algorithm. The E-step computes, for each
point \(i\) and component \(j\), the \emph{responsibility}
\[
  r_{ij} \;=\; \frac{w_j\, \mathcal{N}(x_i;\, \mu_j, \Sigma_j)}
                   {\sum_{l} w_l\, \mathcal{N}(x_i;\, \mu_l, \Sigma_l)} ,
\]
the posterior probability that point \(i\) came from component \(j\). The M-step
re-estimates each component from the points, weighted by responsibility:
\[
  w_j = \frac{1}{n}\sum_i r_{ij}, \qquad
  \mu_j = \frac{\sum_i r_{ij}\, x_i}{\sum_i r_{ij}}, \qquad
  \Sigma_j = \frac{\sum_i r_{ij}\,(x_i - \mu_j)(x_i-\mu_j)^{\mathsf T}}{\sum_i r_{ij}} .
\]
No iteration decreases the likelihood. Mathilda stops when its relative change
falls below \(10^{-10}\), or after 200 iterations. \(k\)-means is the limit of this algorithm in which every
covariance is the same multiple of the identity and shrinks to zero, so each
responsibility becomes \(0\) or \(1\). The number of components is chosen by the
\emph{Bayesian information criterion}\index{Bayesian information criterion}
\[
  \mathrm{BIC} \;=\; q \log n \;-\; 2\log L,
\]
where \(L\) is the maximised likelihood and \(q\) the number of free parameters.
Each extra component adds parameters, and BIC charges \(\log n\) for each one.
Mathilda fits \(K = 1, 2, \dots\) up to at most ten components and keeps the
smallest BIC.

\calloutnote{Theory}{The likelihood of a Gaussian mixture has no maximum. Centre a
component on a single data point and shrink its covariance: its density at that
point, and so the likelihood, goes to infinity. An unguarded BIC search then buys
extra components that each sit on one point. Mathilda adds a floor to each
covariance diagonal. In \mcode{FindClusters} the floor is the square of the median
edge of the minimum spanning tree, on the principle that the data cannot resolve
structure finer than the typical spacing between points. The mixture code in
\texttt{src/ml/gmm.c} also computes responsibilities in log space, with the
log-sum-exp trick, because in several dimensions the densities of far-away points
underflow to zero. A component whose covariance matrix is singular falls back to a
diagonal covariance instead of failing the fit.}

\subsection{Mixtures as densities}
\label{ssec:ml-density}

The same fit, called through \mcode{LearnDistribution}, returns the density itself.
\S\ref{ssec:stat-distributions} used \mcode{LearnDistribution} to fit a single normal
distribution. With \mcode{Method -> "GaussianMixture"} it fits a mixture:

\begin{codepairs}
\pair{machine-learning/mixture/6}{A fitted mixture, printed with its parameters
elided.}
\pair{machine-learning/mixture/7}{The two components have almost equal weight.}
\pair{machine-learning/mixture/8}{The mean of the first is close to the true
\((1.5, 3)\).}
\pair{machine-learning/mixture/9}{\mcode{PDF} on a list of points gives one density
each.}
\pair{machine-learning/mixture/10}{A single Gaussian fitted to the same data.}
\pair{machine-learning/mixture/11}{Between the clouds, the single Gaussian sees
plenty of density. The mixture sees very little.}
\end{codepairs}

The contours of the fitted mixture show the two tilted components:

\plotcell[0.6\linewidth]{machine-learning/gmm-density}

A single Gaussian fitted to this data is centred between the two clouds, in a
region that holds almost no data. The mixture density there is more than ten times
smaller. \mcode{LearnDistribution} also accepts outcomes that are not numbers, with
\mcode{Method -> "ContingencyTable"}, and simply counts them:

\begin{codepairs}
\pair{machine-learning/mixture/12}{A categorical distribution.}
\pair{machine-learning/mixture/13}{The empirical frequencies. An outcome never seen
has probability exactly \(0\).}
\end{codepairs}

\B{SmoothKernelDistribution} completes the set with a kernel density estimate,
used in one of the featured examples below.

\section{Principal components}
\label{sec:ml-pca}

The last unsupervised tool finds the directions in which data vary most.
\emph{Principal component analysis}\index{principal component analysis} (PCA)
rotates the coordinate system so that the first axis points along the direction of
largest variance, the second along the largest remaining direction perpendicular to
the first, and so on. The new coordinates are uncorrelated.

Here is the mathematics. Centre the data, so each column of the \(n\times p\) data
matrix \(X\) has mean zero, and form the sample covariance matrix
\[
  C \;=\; \frac{1}{n-1}\, X^{\mathsf T} X .
\]
\(C\) is symmetric and positive semidefinite, so it has an orthonormal basis of
eigenvectors \(v_1, \dots, v_p\) with eigenvalues
\(\lambda_1 \ge \dots \ge \lambda_p \ge 0\). The variance of the data along a unit
vector \(u\) is \(u^{\mathsf T} C u\), which is largest for \(u = v_1\), where it
equals \(\lambda_1\). The principal components of a row \(x\) are its coordinates
\(v_j^{\mathsf T} x\) in this basis.

\plotcell[0.6\linewidth]{machine-learning/pca-axes}

The figure shows 200 points from a correlated Gaussian, with the two eigenvectors
of \(C\) drawn from the mean, each \(2\sqrt{\lambda_j}\) long in both directions.
The session computes the same quantities:

\begin{codepairs}
\pair{machine-learning/pca/1}{200 correlated points.}
\pair{machine-learning/pca/2}{Their mean and covariance matrix.}
\pair{machine-learning/pca/3}{The eigenvalues and eigenvectors of the covariance
matrix.}
\pair{machine-learning/pca/4}{\mcode{PrincipalComponents}: the data in the rotated
coordinates.}
\pair{machine-learning/pca/5}{The new coordinates have variances
\(\lambda_1, \lambda_2\) and zero covariance.}
\end{codepairs}

Line 5 is the defining property. The variances of the components are the
eigenvalues, and their covariance is zero after \B{Chop} removes rounding at the
\(10^{-16}\) level. The total variance \(3.553 + 1.521 = 4.866 + 0.209\) is
unchanged, because a rotation preserves it. One direction holds \(96\%\) of it.

The first eigenvector is \((-0.847, -0.531)\) according to \B{Eigenvectors}, but
the first principal component of the first point is positive. An eigenvector is
defined only up to sign. \mcode{PrincipalComponents} flips each eigenvector so that its
largest entry is positive, so its output does not depend on which eigensolver was
used.

\calloutnote{Theory}{PCA is the \emph{singular value decomposition} of the
centred data. If \(X = U S V^{\mathsf T}\), then
\(X^{\mathsf T}X = V S^2 V^{\mathsf T}\), so the right singular vectors are the
eigenvectors of \(C\) and \(\lambda_j = s_j^2/(n-1)\). Many libraries compute PCA
this way, because it avoids forming \(X^{\mathsf T}X\) and so avoids squaring the
condition number. Mathilda's \texttt{ml\_pca} (\texttt{src/ml/pca.c}) instead forms
the \(p \times p\) covariance matrix and diagonalises it, with LAPACK's
\texttt{dsyev} when it is linked and a cyclic Jacobi method otherwise. For a
modest number of features the two give the same answer, as the next two lines
show.}

\begin{codepairs}
\pair{machine-learning/pca/6}{The right singular vectors of the centred data, the
same as the eigenvectors in line 3.}
\pair{machine-learning/pca/7}{\(s_j^2/(n-1)\) gives the eigenvalues again.}
\end{codepairs}

\subsection{Reducing dimension}
\label{ssec:ml-reduce}

Keeping only the first \(k\) principal components reduces each row from \(p\)
numbers to \(k\), with the least loss of variance possible for a linear map.
\mcode{DimensionReduce} returns the reduced data. \mcode{DimensionReduction} returns a
reusable \mcode{DimensionReducerFunction} that applies the same projection to new
points:

\begin{codepairs}
\pair{machine-learning/pca/8}{The first component of the first three points, as in
line 4.}
\pair{machine-learning/pca/9}{A reducer from two dimensions to one.}
\pair{machine-learning/pca/10}{It projects a new point, or a matrix of points. The
training mean projects to (almost) zero.}
\pair{machine-learning/pca/11}{\mcode{Method -> "Correlation"} standardises each
column first.}
\pair{machine-learning/pca/12}{Points on a line have all their variance in one
component.}
\end{codepairs}

The reducer subtracts the \emph{training} mean before it projects, which makes
projections of different batches comparable. The correlation method is the right
choice when the columns have different units. Without it, the column with the
largest numbers dominates the first component for no statistical reason.
\mcode{DimensionReduce} also offers \mcode{"LatentSemanticAnalysis"}, which skips the
centring (a truncated SVD, suited to sparse non-negative data such as word counts),
and \mcode{"MultidimensionalScaling"}, which works from the matrix of pairwise
distances and is limited to 2000 rows because that matrix is \(n \times n\). Only
PCA is available as a reusable reducer.

\section{Performance and limits}
\label{sec:ml-performance}

The cost of each method follows from its algorithm. Table~\ref{tab:ml-cost} gives
the order of the work for \(n\) training points, \(p\) features, and \(K\)
classes. It is read from the implementations in \texttt{src/ml/}, not measured.

\begin{table}[ht]
\centering
\small
\begin{tabular}{@{}lll@{}}
\toprule
Method & training & one prediction \\
\midrule
nearest neighbours & copy the data, \(O(np)\) & scan the data, \(O(np)\) \\
naive Bayes & one pass, \(O(np)\) & \(O(Kp)\) \\
logistic regression & per step \(O(np^2)\), times \(K\) for one-vs-rest & \(O(Kp)\) \\
decision tree & per level about \(O(pn\log n)\) & depth of the tree \\
random forest & 50 trees on bootstrap samples & 50 tree walks \\
least squares & \(O(np^2 + p^3)\) & \(O(p)\) \\
\(k\)-means (\mcode{FindClusters}) & \(O(nkp)\) per iteration & --- \\
Gaussian mixture & \(O(nKp^2)\) per EM step, for each \(K\) tried & --- \\
PCA & \(O(np^2 + p^3)\) & \(O(kp)\) \\
\bottomrule
\end{tabular}
\caption{Order of the work for each method.}
\label{tab:ml-cost}
\end{table}

The repository has measurements for the clustering code. The reference page for
\mcode{FindClusters} (\texttt{docs/spec/builtins/lists-and-iteration.md}) records that
the spanning-tree methods on machine-precision points take \(6.8\) ms for 2000
points, \(0.17\) s for 10\,000, and \(0.81\) s for 20\,000, and are capped at
20\,000 points because the cost grows quadratically. Exact (rational) points and
strings are capped at 2000. On a line there is no cap, and \(10^6\) numbers take
about \(2.3\) s. \mcode{"KMeans"} on 20\,000 points in ten dimensions with
\(k = 3\) takes \(0.60\) s. The machine-learning reference page
(\texttt{docs/spec/builtins/machine-learning.md}) records the 2000-row cap on
multidimensional scaling.

There are also benchmarks of machine-learning algorithms written directly in
Mathilda's array language instead of calling a builtin. The third benchmark sweep
in \texttt{docs/design/performance.md} (\S9) and experiment 10 in
\texttt{docs/experiments/} time a \(k\)-means of 100\,000 points in eight
dimensions with \(k = 16\) for 20 iterations. Mathilda takes \(2.26\) s,
Mathematica~14.0 \(7.33\) s, and NumPy \(0.77\) s. A logistic regression on a
\(200\,000 \times 32\) matrix with 100 gradient-descent steps takes \(2.71\) s in
Mathilda, \(2.15\) s in Mathematica, and \(0.84\) s in NumPy. Those two programs
spend their time in \B{Dot} and element-wise arithmetic on packed arrays
(\S\ref{sec:ds-packed}). The same report attributes the remaining gap to NumPy to
several passes over the data where a hand-fused loop would make one.

\calloutnote{Performance}{A fitted model is kept as ordinary lists of numbers.
Every application reads the parameters it needs back out of those lists, and a
nearest-neighbour or forest prediction reads a lot of them: the whole training set,
or fifty trees. This keeps models printable, comparable with \mcode{SameQ}, and
editable with \mcode{Part}, at the cost of a slower prediction. There are no
benchmarks of \mcode{Classify} or \mcode{Predict} in the repository, so this chapter
quotes no timings for them.}

\subsection*{What is missing}

Compared with Mathematica~15, the machine-learning subsystem is small, and it
does not hide the gaps: an unsupported method, option, or data layout returns the
call unevaluated. The main differences are these.

\begin{itemize}
\item Features must be real numbers. There is no automatic handling of
  categorical, text, image, date, or missing values, and no
  \mcode{FeatureExtraction}.
\item \mcode{Classify} accepts only a list of rules, and a model applies to one point
  at a time.
\item \mcode{Classify} has five methods. Mathematica also has support vector machines,
  gradient-boosted trees, neural networks, and Markov models. \mcode{Predict} has only
  least squares and nearest neighbours.
\item There is no automatic method selection, no hyperparameter tuning, no
  \mcode{ValidationSet}, and no class weights. The forest always has fifty trees,
  and trees are never pruned.
\item There is no \mcode{ClassifierMeasurements}, \mcode{PredictorMeasurements},
  or \mcode{ClusterClassify}. This chapter builds accuracy, the confusion matrix,
  cross-validation, and cluster assignment from a few lines each.
\item A model answers only a few properties: \mcode{"Classes"},
  \mcode{"Method"}, \mcode{"FeatureCount"}, \mcode{"NeighborCount"}, and for a
  predictor \mcode{"Coefficients"} and \mcode{"TrainingData"}. Everything else is
  in the payload, which you can read with \mcode{Part}, as this chapter does.
\end{itemize}

\section{Featured examples}
\label{sec:ml-featured}

The examples in this section are short tasks on real data. Each is a few lines of
input, usually with a picture. The iris examples share one session, so later
examples use names defined in earlier ones.

\subsection*{Classify iris species}

In 1936 R.~A. Fisher published measurements of 150 iris flowers from three
species: sepal length, sepal width, petal length, and petal width, in
centimetres. The session types in the first fifteen flowers of each species,
trains on ten of each, and tests on the other five.\index{iris data}

\begin{codepairs}
\pair{machine-learning/iris/1}{The first fifteen flowers are \emph{Iris setosa}.}
\pair{machine-learning/iris/2}{The next fifteen are \emph{Iris versicolor}.}
\pair{machine-learning/iris/3}{The last fifteen are \emph{Iris virginica}.}
\pair{machine-learning/iris/4}{}
\pair{machine-learning/iris/5}{Ten of each species for training, five of each for
testing.}
\pair{machine-learning/iris/6}{A logistic-regression classifier.}
\pair{machine-learning/iris/7}{Predictions for the fifteen test flowers.}
\pair{machine-learning/iris/8}{All fifteen are right.}
\pair{machine-learning/iris/9}{A flower between versicolor and virginica gets a
mixed answer.}
\pair{machine-learning/iris/10}{The other deterministic methods get 14, 15, and
15 right.}
\end{codepairs}

Setosa is easy, because its petals are much smaller. Versicolor and virginica
overlap. In the petal plane, five nearest neighbours draw these regions:

\plotcell[0.6\linewidth]{machine-learning/iris-regions}

\subsection*{Read the rules a tree learned}

A decision tree is a classifier you can read. The tree for the iris training set
has five nodes:

\begin{codepairs}
\pair{machine-learning/iris/11}{Split on feature 2 (petal length) at \(2.5\), then
on feature 3 (petal width) at \(1.65\).}
\pair{machine-learning/iris/12}{The class counts at each node: each leaf holds ten
flowers of one species.}
\end{codepairs}

As rules: if the petal length is at most \(2.5\) cm, the flower is a setosa.
Otherwise, if the petal width is at most \(1.65\) cm, it is a versicolor, and if
not, a virginica. Feature numbers start at zero, so feature 2 is the third column.
The threshold \(2.5\) lies halfway between the longest setosa petal in the
training set (\(1.7\)) and the shortest versicolor petal (\(3.3\)).

\subsection*{Compress measurements with PCA}

Four measurements per flower are correlated, because big flowers are big in every
dimension. Principal components show how much of the variation is really
independent:

\begin{codepairs}
\pair{machine-learning/iris/13}{Principal components of all 45 flowers.}
\pair{machine-learning/iris/14}{The fraction of the total variance in each
component.}
\pair{machine-learning/iris/15}{A reusable reducer projects a new flower to two
numbers.}
\end{codepairs}

The first component alone holds \(92\%\) of the variance, and the first two hold
\(98\%\). The projection onto those two keeps the species apart:

\plotcell[0.6\linewidth]{machine-learning/iris-pca}

\subsection*{Discover species without labels}

Suppose the species names were lost. Can clustering recover them?

\begin{codepairs}
\pair{machine-learning/iris/16}{\(k\)-means with three clusters, cross-tabulated
against the species (columns: setosa, versicolor, virginica).}
\end{codepairs}

The first cluster is exactly the fifteen setosa flowers. The other two clusters
each mix versicolor and virginica, and they agree with the species on 24 of the
30 flowers. Here are the clusters, in the petal plane:

\plotcell[0.6\linewidth]{machine-learning/iris-clusters}

\subsection*{See two modes in one measurement}

The petal lengths alone separate setosa from the rest. A kernel density estimate
(\mcode{SmoothKernelDistribution}) makes that visible without any labels. It places a
normal density on each observation and averages them, with a bandwidth chosen by
the normal-reference rule \(h = \hat\sigma\,(4/(3n))^{1/5}\).

\plotcell[0.6\linewidth]{machine-learning/iris-bimodal}

The printed bandwidth is \(0.89\) cm. The left mode is setosa. The right one
merges versicolor and virginica, whose petal lengths overlap. The
normal-reference rule is designed for data close to a single normal distribution,
so on multimodal data it tends to oversmooth.

\subsection*{Predict fuel economy}

The 1974 \emph{Motor Trend} road tests (the \texttt{mtcars} data set) give the
weight, horsepower, and fuel economy of 32 cars. Heavier cars use more fuel:

\begin{codepairs}
\pair{machine-learning/mileage/1}{Load the names of the 32 cars.}
\pair{machine-learning/mileage/2}{Their weights are given in thousands of pounds.}
\pair{machine-learning/mileage/3}{Gross horsepower follows.}
\pair{machine-learning/mileage/4}{Fuel economy is measured in miles per gallon.}
\pair{machine-learning/mileage/5}{About \(5.3\) miles per gallon lost per 1000
pounds.}
\pair{machine-learning/mileage/6}{Predicted economy of a 3000-pound car.}
\end{codepairs}

The coefficients \(37.285\) and \(-5.344\) are the standard textbook result for
this data set.

\plotcell[0.6\linewidth]{machine-learning/mileage-line}

\subsection*{Weigh two predictors at once}

Weight and horsepower are correlated: heavy cars tend to have big engines. A
regression on both separates their effects:

\begin{codepairs}
\pair{machine-learning/mileage/7}{Intercept, then the coefficients of weight and
horsepower.}
\pair{machine-learning/mileage/8}{A 3000-pound car with 150 horsepower.}
\end{codepairs}

With horsepower in the model, the weight coefficient drops from \(-5.34\) to
\(-3.88\). Part of what looked like the effect of weight was the effect of the
engines that heavy cars carry.

\subsection*{Flag unusual cars}

A density estimate can flag unusual cases. Fit a two-dimensional normal
distribution to (weight, economy) and list the cars where the fitted density is
lowest:

\begin{codepairs}
\pair{machine-learning/mileage/9}{A bivariate normal fitted to weight and
economy.}
\pair{machine-learning/mileage/10}{The three cars with the lowest density.}
\end{codepairs}

The Chrysler Imperial is heavy but more economical than its weight suggests. The
Toyota Corolla and the Fiat 128 are light and unusually economical even for their
weight. All three lie far from the regression line in the previous figure.

\subsection*{Detect spam messages}

Mathilda's classifiers take numbers, so text must first become features. This
example uses three: the number of exclamation marks, the fraction of capital
letters, and the number of words from a short list of money words. The training
set is twelve short messages, typed in:

\begin{codepairs}
\pair{machine-learning/spam/1}{Three numerical features of a message.}
\pair{machine-learning/spam/2}{Six of the messages are spam.}
\pair{machine-learning/spam/3}{The other six are genuine.}
\pair{machine-learning/spam/4}{A naive Bayes filter.}
\pair{machine-learning/spam/5}{The features of a new message.}
\pair{machine-learning/spam/6}{Three new messages, classified.}
\pair{machine-learning/spam/7}{One innocent use of ``free'' is enough to condemn a
message.}
\end{codepairs}

The last line shows the limits of twelve examples. No genuine message in the
training set contains a money word, so within that class the feature has zero
variance, and naive Bayes treats any money word as proof of spam (see the pitfall
in \S\ref{sec:ml-bayes}). A real filter needs thousands of messages and many
more features, but the pipeline is the same: text to numbers, numbers to a
classifier.

\subsection*{Recognise noisy digits}

Each digit is drawn on a grid three pixels wide and five tall, and flattened into
a vector of fifteen numbers. The training and test sets are copies with Gaussian
noise added to every pixel. One noisy copy of each digit:

\plotcell[0.7\linewidth]{machine-learning/digits-noisy}

\begin{codepairs}
\pair{machine-learning/digits/1}{The ten clean digits.}
\pair{machine-learning/digits/2}{Noise with standard deviation \(0.3\), clipped to
\([0, 1]\).}
\pair{machine-learning/digits/3}{Twenty noisy copies of each digit for
training.}
\pair{machine-learning/digits/4}{Thirty more of each for testing.}
\pair{machine-learning/digits/5}{\(88\%\) of the test digits are read
correctly.}
\pair{machine-learning/digits/6}{A noisy 8.}
\pair{machine-learning/digits/7}{The most common mistakes, as \{\{true, read\},
count\}.}
\end{codepairs}

The most common mistake is reading a 0 as an 8. The two digits differ in a single
pixel, the centre of the grid, so noise on that one pixel can turn one into the
other. The next is 9 read as 5, which also differ in one pixel. Nearest-neighbour
classification of images is a real method, and on the classic MNIST digits it does
well with enough training data. Its weakness is also visible here: pixel
distance does not know which pixels matter.

\subsection*{Smooth a noisy signal}

Nearest-neighbour regression is also a smoother. Train it on noisy samples of a
sine wave, and it returns a local average:

\plotcell[0.6\linewidth]{machine-learning/smooth}

The orange curve is the five-neighbour predictor, and the blue curve is the true
sine. The predictor follows the signal well in the middle of the range. At the two
ends its neighbours all lie on one side, so it flattens out. This bias at the
boundary is shared by every local-averaging smoother.

\subsection*{Where this connects}

Machine learning in Mathilda is built from parts that earlier chapters cover. The
training data are lists and rules (Chapter~\ref{ch:datastructures}), and a fitted
model is an ordinary expression that \mcode{Part}, \mcode{FullForm}, and \mcode{SameQ} work
on. The random data come from the seeded generators of
\S\ref{ssec:stat-distributions}, so every example here can be reproduced exactly.
Least squares and principal components are linear algebra (\S\ref{sec:linalg}):
the session checked \mcode{Predict} against \mcode{LeastSquares} and
\mcode{PrincipalComponents} against \B{SingularValueDecomposition}. Single-linkage
clustering is a minimum spanning tree, one of the graph objects of
Chapter~\ref{ch:graphs}, and the array benchmarks quoted in
\S\ref{sec:ml-performance} run on the packed arrays of \S\ref{sec:ds-packed}. For
the full options of any function in this chapter, type \texttt{?Name} at the
prompt.
